% EGA V, English translation by P. Blass, 1987 (cote 24): archive cote 24, site page 76 (PDF page 77).
% typescript. Draft compared: none.
% AI-transcribed by Claude Opus 5.5 (claude-opus-5-5, adaptive thinking, effort high), Message Batches API, batch ega24b (msgbatch_01RGJJeHdR96U16ZWFuUjtDx), 2026-11-13; not yet checked by eye.
\scanpage{24/76}
\mspage{-3-}

fondation d'un mouvement international et interprofessionnel
``Survivre'', et à consacrer aux questions liées à la survie une
part importante de mon énergie. Dans cette optique, la seule va-%
\margin{(}leur de mon apport comme mathématicien est de me permettre aujourd'
hui, grâce à l'estime professionnelle et personnelle acquise parmi
mes collègues, de donner plus de force à mon témoignage et à mon
action en faveur d'une stricte subordination de toutes nos acti-%
vités, y compris nos activités de scientifiques, aux impératifs
\margin{(}de la survie, et à la promotion d'un ordre stable et humain sur
notre planète, sans lequel la survie de notre espèce ne serait
ni possible, ni désirable.

A. GROTHENDIECK.

PRINCIPALES PUBLICATIONS
------------------------

\uline{Espaces Vectoriels Topologiques}

[1] Critères de compacité dans les espaces fonctionnels géné-%
raux, Amer. J. 74 (1952), p. 168-186.

[2] Sur certains espaces de fonctions holomorphes, J. Crelle
192 (1953), p. 35-64 et 77-95.

[3] Espaces Vectoriels Topologiques, Notes polyc., Sa\~o Paulo
(1954), 240 p.

[4] Sur les espaces (F) et (DF), Summa Bras. 3 (1954),
p. 57-123.

[5] Produits tensoriels topologiques et espaces nucléaires,
Mem. AMS, n° 16 (1955), 329 p.

[6] Résumé de la théorie métrique des produits tensoriels
topologiques, Bull. Sa\~o Paulo 8 (1953), p. 1-79.

[7] La théorie de Fredholm, Bull. SMF, 84 (1956), p. 319-384.

\uline{Topologie et algèbre homologique}

[8] Sur quelques points d'algèbre homologique, Tohoku M.J.,
9 (1957), p. 119-221.

[9] Théorèmes de finitude pour la cohomologie des fasceaux,
Bull. SMF, 84 (1956), p. 1-7.

\uline{Géométrie analytique}

[10] Sur la classification des fibrés holomorphes sur la
sphère de Riemann, Amer. J., 79 (1957), p. 121-138.

[11] Techniques de construction en géométrie analytique,
Sem. H. Cartan, 13 (1960/61), exposés 7 à 17.

...

\margin{76}
